\chapter{Теория}

\section{Типы обрабоки коцнов витой пары}


\begin{figure}[H] 
    \centering
    \includegraphics[width=0.75\textwidth]{wp.png}
    \caption{Типы обработки}
\end{figure}


Сетевые кабели состоят из четырех пар медных проводов различных цветов. Каждая пара
проводов скручена с разным шагом, что способствует уменьшению электромагнитных помех.
На рисунке 7 показаны различные варианты обработки концов витой пары. \\

В сетях Ethernet 10/100BASE-T используются только две пары проводов - оранжевая и
зеленая. Две оставшиеся пары, синяя и коричневая, используются для других приложений
Ethernet и телефонных соединений. \\

Также существуют два способа обработки концов кабелей - прямой и перекрестный. \\

Прямой способ (straight through cable) предполагает обработку обоих концов кабеля по
типу A или B. Прямой кабель используется для подключения разных уровней устройств,
например, компьютера к коммутатору или маршрутизатора. \\

Перекрестный способ (crossover cable) предполагает обработку одного конца кабеля по
типу A, а другого - по типу B. Перекрестный кабель используется для подключения устройств
на одном уровне, например, компьютера к компьютеру, коммутатора к коммутатору или
маршрутизатора к маршрутизатору. \\

На данный момент прямой кабель с обработкой по типу B является более популярным и
широко используется, так как современные сетевые устройства не требуют строгой
совместимости. Кабели типа A и перекрестные кабели почти не используются.

\section{Последовательность оконечивания кабеля}

\begin{enumerate}
    \item Снять верхнюю оболочку кабеля.
    \item Скрученные провода раскрутить.
    \item Расположить провода в нужном порядке (как на картинке выше)
    \item Взять коннектор и вставить в него провода в нужном порядке.
    \item Обжать коннектор с помощью кримпера.
\end{enumerate}

\section{Проверка работоспособности кабеля}

Для проверки работоспособности существуют спецаильные устройства - тестеры. \\

\begin{figure}[H] 
    \centering
    \includegraphics[width=0.75\textwidth]{tester.png}
    \caption{Тестер для кабеля}
\end{figure}

Тестер будет последовательно проверять сигнал по всем проводам. Если провод
исправен, то зеленый индикатор загорится напротив его порядкового номера. \\

При тестировании прямого кабеля, индикаторы на левой и правой стороне тестера должны
гореть синхронно и по порядку (1 - 1, 2 - 2, 3 - 3... 8 - 8). Если все индикаторы горят зеленым и
последовательность правильная, то кабель считается исправным. \\

При тестировании перекрестного кабеля, лампочки на левой и правой стороне тестера
должны гореть синхронно в следующей последовательности (1-3, 2-6, 3-1, 4-4, 5-5, 6-2, 7-7, 8-
8). Если все лампочки горят зеленым и последовательность правильная, то кабель считается
исправным. \\

В случае, если лампочка не горит или горит красным, рекомендуется 
усилить обжатие кабеля. Затем повторно проверьте кабель. Если это не помогло, то следует
обрезать коннектор и снова обжать кабель.
\endinput